This project sought to create a computational model that characterised vibrations propagating through an intraosseous transcutaneous amputation prosthesis (ITAP).
It aimed to set the path for further studies, wherein we could model bone-anchoring, whether the device was press-fit or cemented, etc.
It did so by creating a computational model using Abaqus, and carrying out both modal analysis and the application of a concentranted point-force.
Sensors were placed along the device, and their displacements were monitored.

\paragraph{Modal analysis}
For the particular configuration of ITAP used, we observed natural frequencies of vibration at about \SI{930}{Hz} and \SI{3930}{Hz}, informing the maximum range of frequencies to be explored in our study, and confirming that it is well-behaved in the clinically relevant frequencies.

\paragraph{Direct control of amplitude}
Given the nature of our model, we could directly control the amplitude of vibrations by a) controlling the amplitude of the input and b) shifting the position of the input.
As the input position approaches the fulcrum point, we saw a decline in the maximum amplitude across all three sensors in the device.
Interestingly, each sensor had a different decreasing rate, which suggested we could directly control the relative displacement between each part by changing the position of the actuator. 

\paragraph{Frequency analysis}
We observed a linear response up to about \SI{700}{Hz} with a decline in the maximum displacement of \SI{60}{\deci\bel\per decade}. Particularly above \SI{400}{Hz} we also observe a frequency-dependent increase in the relative displacement of the stem with respect to the spigot.

\paragraph{Time analysis}
For frequencies at and below \SI{400}{Hz} we observe a time-invariant behaviour.
Beyond it, we see a frequency-dependent amplitude loss over time, which becomes both larger and faster with a frequency increase --- culminating at 30\% for \SI{850}{Hz}.